\section{Evaluación y resultados}

\subsection{Por qué se construyó un banco de pruebas}

A diferencia de un error de programación, que se manifiesta con
una excepción, los errores producidos por el modelo de lenguaje son silenciosos: 
el sistema puede generar una respuesta errónea sin que se detecte de inmediato.

Durante el desarrollo se detectó un caso que ilustra el problema. Ante la
pregunta ``¿he gastado más este año que el año pasado?'', el modelo generaba una
consulta con \texttt{strftime('\%Y-\%m')} en lugar de \texttt{strftime('\%Y')}.
La consulta comparaba el mes actual contra el mes anterior, pero la respuesta
los presentaba como años. El asistente respondía que se había gastado más este
año cuando la realidad era 17.870 € frente a 26.643 €: la conclusión estaba
invertida y se expresaba con seguridad. Revisando la interfaz, la respuesta
resultaba convincente.

Por este motivo se desarrolló \texttt{tests/evaluar.py}, un banco de preguntas
que se ejecuta contra el agente real y mide su comportamiento de forma
sistemática.

\subsection{Qué se mide y por qué se separa}

El banco contiene 44 preguntas etiquetadas, distribuidas entre los tipos de
consulta que menciona el enunciado (gastos, comparativas y suscripciones) y
otros añadidos durante el desarrollo: seguimiento multiturno, operaciones de
Bizum, saldo, ingresos, preguntas ambiguas, límites del asistente y previsión.

Se puntúan cuatro indicadores por separado, porque responden a fallos de
naturaleza distinta:

\begin{itemize}
    \item \textbf{Elección de herramienta.} Si el usuario pregunta por sus 
    suscripciones y el modelo decide escribir una consulta SQL manual en lugar 
    de invocar \texttt{analizar\_suscripciones}, se considera un fallo aunque 
    la cifra final resulte ser correcta. Este indicador se evalúa mediante un 
    espía sobre la función \texttt{ejecutar\_tool}, permitiendo registrar 
    también las peticiones que el \emph{backend} resuelve de forma directa 
    (como el saldo) y que no dejan rastro en el historial del modelo.

    \item \textbf{Respuesta final.} Que la consulta sea correcta no garantiza
    que el modelo traslade bien la cifra al texto. Es lo único que oye el
    usuario, y se comprueba que los valores esperados aparezcan con formato
    español.

    \item \textbf{Refuerzo visual.} Se puntúa en los dos sentidos: no mostrar
    nada cuando hay varios valores comparables es un fallo, y mostrar algo
    cuando la respuesta es un único dato también. Una tabla cuenta igual que un
    gráfico, porque la elección entre ambos corresponde al modelo y las dos son
    respuestas válidas.

    \item \textbf{Validez de la especificación.} Que llegue un gráfico no
    significa que se vea. Una especificación sin el campo \texttt{mark} supera
    la validación del backend y falla después en el renderizado, sin dejar
    rastro. Se puntúa por separado del indicador anterior porque son fallos de
    cosas distintas: uno mide al modelo decidiendo y el otro al modelo
    redactando.
\end{itemize}

Adicionalmente, medimos un quinto indicador: ejecutamos la consulta SQL generada por el agente en paralelo a 
una consulta SQL escrita a mano (la referencia), y comparamos los resultados. 
No se puntúa porque el modelo puede interpretar la pregunta de forma diferente 
pero válida. Además, nunca comparamos el texto en crudo de la consulta SQL, 
ya que existen multitud de sentencias válidas para obtener un mismo resultado.

Por último, evaluamos la latencia en dos tramos. El tramo más relevante mide 
el tiempo transcurrido hasta emitir el evento \texttt{fin\_respuesta}, que marca 
el momento en el que el texto de la respuesta está completo. Esta métrica 
representa una cota superior de la espera percibida por el usuario, ya que la 
síntesis de voz arranca antes (con la primera frase).

\subsection{Resultados}

Las métricas se obtuvieron realizando tres repeticiones independientes de las 44 
preguntas, utilizando el modelo local \texttt{qwen3:8b} y regenerando la base de 
datos sintética antes de cada evaluación.

\begin{table}[h]
\centering
\begin{tabular}{lr}
\toprule
\textbf{Indicador} & \textbf{Resultado} \\
\midrule
Elección de herramienta            & 97,6 \% \\
Respuesta final correcta           & 88,9 \% \\
Refuerzo visual cuando corresponde & 100 \% \\
Especificaciones válidas           & 100 \% \\
Latencia mediana hasta la voz      & 3,3 s \\
\bottomrule
\end{tabular}
\caption{Resultados sobre 44 preguntas y 3 repeticiones.}
\end{table}

\subsection{Decisiones de diseño tomadas a partir de la medición}

El banco de pruebas no se utilizó únicamente para verificar el resultado final.
En varios casos, la medición contradijo una decisión previa y motivó un cambio
de arquitectura.

\subsubsection{La decisión de mostrar un gráfico}

Inicialmente, la decisión de acompañar una respuesta con un gráfico se dejaba al
modelo mediante instrucciones en el \textit{system prompt}. La medición mostró
que esa decisión no era estable: variaba ante cualquier modificación del
\textit{prompt}, incluso ante cambios que no trataban sobre gráficos.

\begin{table}[h]
\centering
\begin{tabular}{lcccc}
\toprule
\textbf{Pregunta} & \textbf{v1} & \textbf{v2} & \textbf{v3} & \textbf{v4} \\
\midrule
Ranking de comercios          & 0/3 & 2/3 & 0/3 & 0/3 \\
Comparación mensual           & 0/3 & 0/3 & 0/3 & 0/3 \\
Comparación por categorías    & 0/3 & 2/3 & 2/3 & 0/3 \\
Serie de seis meses           & 0/3 & 3/3 & 3/3 & 0/3 \\
Comparación interanual        & 0/3 & 0/3 & 0/3 & 0/3 \\
Evolución mensual             & 3/3 & 3/3 & 3/3 & 3/3 \\
\bottomrule
\end{tabular}
\caption{Veces que cada pregunta produce un gráfico, sobre cuatro versiones del
\textit{system prompt}. Entre la tercera y la cuarta solo cambiaron dos líneas
relativas al tratamiento de fechas.}
\end{table}

A partir de este resultado se reorganizó la responsabilidad. La decisión de
\textit{si} existe algo que representar pasó al backend, mediante una regla
determinista sobre la forma del resultado: varias filas con alguna columna
numérica, o una sola fila con dos o más cifras comparables. La decisión de
\textit{qué} representación utilizar (tabla o gráfico), con qué marca, qué ejes
y con qué justificación, permanece en el modelo, en una llamada dedicada de
tarea única.

El indicador pasó del 16,7 \% al 100 \%, sin variación en la precisión de las
respuestas ni en la latencia percibida. La Figura~\ref{fig:resultados} recoge
ese cambio junto al de la latencia, que es la otra decisión que se tomó a partir
de una medición.

\begin{figure}[htbp]
\centering
\begin{tikzpicture}[
    font=\small,
    barra/.style={draw=#1, fill=#1!25, thick},
    barra/.default=gray,
    val/.style={font=\footnotesize\bfseries, text=#1, anchor=south, inner sep=2pt},
    val/.default=gray!60!black,
    pie/.style={font=\scriptsize, text=gray!55!black, align=center, anchor=north,
                inner sep=3pt, text width=2.4cm},
    titulo/.style={font=\footnotesize\bfseries, text=subsectionblue, anchor=south,
                   align=center, text width=6.4cm},
]

% ================= Panel 1: refuerzo visual (mayor es mejor) =============
% escala: 3 cm = 100 %
\node[titulo] at (2.3,3.6) {Preguntas que reciben el elemento visual que les corresponde};

\draw[barra=ugrred]         (0.5,0) rectangle (1.6,0.50);
\node[val=ugrred]        at (1.05,0.50) {16,7 \%};
\node[pie]               at (1.05,-0.1) {lo decide el modelo};

\draw[barra=subsectionblue] (3.0,0) rectangle (4.1,3.00);
\node[val=subsectionblue] at (3.55,3.00) {100 \%};
\node[pie]                at (3.55,-0.1) {regla determinista en el backend};

\draw[gray!55] (0.2,0) -- (4.4,0);

% ============ Panel 2: latencia hasta la voz (menor es mejor) ============
% escala: 3 cm = 12,4 s; solo preguntas que generan un elemento visual
\node[titulo] at (10.0,3.6) {Tiempo hasta la respuesta hablada en preguntas con elemento visual};

\draw[barra=ugrred]         (8.2,0) rectangle (9.3,3.00);
\node[val=ugrred]        at (8.75,3.00) {12,4 s};
\node[pie]               at (8.75,-0.1) {la herramienta visual se llama dentro del bucle};

\draw[barra=subsectionblue] (10.7,0) rectangle (11.8,1.62);
\node[val=subsectionblue] at (11.25,1.62) {6,7 s};
\node[pie]                at (11.25,-0.1) {llamada dedicada tras \texttt{fin\_respuesta}};

\draw[gray!55] (7.9,0) -- (12.1,0);

\end{tikzpicture}
\caption{Las dos decisiones que más cambiaron al medirlas. A la izquierda,
mover al \emph{backend} la decisión de \textit{si} mostrar algo; a la derecha,
sacar la generación del elemento visual fuera del camino crítico. La segunda
medición corresponde solo a las preguntas que acaban generando una tabla o un
gráfico, no a la mediana de las 44.}
\label{fig:resultados}
\end{figure}


\subsubsection{El tamaño del \textit{system prompt}}

Durante el desarrollo, el \textit{system prompt} fue creciendo al incorporar
nuevas reglas. La medición mostró que ese crecimiento degradaba reglas
anteriores. La causa era la duplicación: las mismas seis herramientas se describían tanto en
el \textit{prompt} como en sus esquemas JSON, y ambos textos se envían en cada
petición. Se eliminó la parte duplicada (la que indica cuándo llamar a cada
herramienta, que el esquema ya expresa y en mejor posición, junto a la
definición de la función) y se conservó la que el esquema no puede expresar:
cómo interpretar el resultado devuelto.

\begin{table}[h]
\centering
\begin{tabular}{lrr}
\toprule
 & \textbf{Antes} & \textbf{Después} \\
\midrule
Contexto fijo por petición   & 3.822 tokens & 3.208 tokens \\
Porcentaje de la ventana     & 47 \%        & 39 \% \\
Elección de herramienta      & 97,6 \%      & 97,6 \% \\
Respuesta final correcta     & 85,7 \%      & 90,5 \% \\
\bottomrule
\end{tabular}
\caption{Efecto de eliminar las instrucciones duplicadas.}
\end{table}

Se verificó una por una que ninguna de las diecisiete reglas obtenidas mediante
medición se hubiera perdido en el proceso. Se recuperaron dos preguntas, una de
ellas registrada desde el inicio como limitación asumida del modelo: no lo era.

\subsection{Validación de los componentes deterministas}

Los componentes que no dependen del modelo se validan por separado, con pruebas
que no requieren el LLM y que por tanto son deterministas y rápidas.

\subsubsection{Detección de pagos recurrentes}

Los umbrales del detector se ajustaron observando una única base de datos, lo
cual no demuestra que generalicen. Para comprobarlo se regenera el histórico con
500 semillas distintas y se contrasta el resultado con la verdad conocida por
construcción: 0 falsos negativos, 1 falso positivo (0,2 \% de los históricos
generados), ninguna cadencia mal clasificada y el seguro anual detectado en las
500 semillas.

Conviene señalar el límite de esta validación. El generador asigna a los pagos
recurrentes un día fijo de cada mes, mientras que los gastos no recurrentes caen
en un día aleatorio. La prueba demuestra por tanto que el detector distingue un
cargo de día fijo de uno de día aleatorio, y que no clasifica como recurrente un
gasto habitual: una compra semanal, incluso perfectamente regular, no se detecta
como suscripción, porque el clasificador solo admite cadencias reconocidas. No
cubre, en cambio, el desplazamiento del día de cobro habitual en las
domiciliaciones reales, que se cargan el siguiente día hábil. Medido
posteriormente, el detector reconoce la cadencia con desplazamientos de hasta
dos días, y deja de reconocerla a partir de tres.

\subsubsection{Proyección de gasto}

El estimador de la proyección a fin de mes se seleccionó comparando tres
alternativas sobre los 24 meses completos del histórico, midiendo el error de
cada una según el día del mes en que se realiza la estimación.

\begin{table}[h]
\centering
\begin{tabular}{lrrrr}
\toprule
\textbf{Método} & \textbf{Día 5} & \textbf{Día 10} & \textbf{Día 15} & \textbf{Día 20} \\
\midrule
Proporción directa            & 174,5 \% & 73,3 \% & 46,1 \% & 25,9 \% \\
Pagos fijos por separado      & 31,7 \%  & 18,4 \% & 13,6 \% & 10,6 \% \\
Combinado con media histórica & 8,7 \%   & 7,6 \%  & 8,4 \%  & 7,7 \% \\
\bottomrule
\end{tabular}
\caption{Error medio de cada estimador según el día del mes.}
\end{table}

La proporción directa resulta inservible porque el alquiler y los recibos
domiciliados se cargan a principios de mes: el día 3, con 890,34 € gastados de
los cuales 844,80 € corresponden a pagos fijos, una regla de tres proyectaría
8.903 € para el mes completo.

El estimador adoptado no extrapola los pagos fijos, cuyo coste mensual ya conoce
el detector de recurrencia, y combina el ritmo de gasto variable observado con
la media histórica, ponderando según lo avanzado que esté el mes. El error se
mantiene estable en torno al 8 \% durante todo el mes, lo que permite ofrecer
una previsión útil desde los primeros días.

Cada proyección se acompaña de su margen de error, calculado aplicando el mismo
estimador a cada mes pasado y midiendo la desviación. El margen varía
considerablemente según la categoría: 0 \% en alquiler o gimnasio, que son pagos
fijos conocidos; en torno al 20 \% en supermercado o gasolina; y superior al
100 \% en los Bizums enviados, que son impredecibles por naturaleza. Ofrecer una
cifra única para todas las categorías habría resultado igual de creíble y
sustancialmente menos honesto.

\subsubsection{Flujo de operaciones}

El envío de Bizums se valida mediante 35 comprobaciones que recorren el flujo
completo sin intervención del modelo: reconocimiento de la petición, validación
del importe, resolución del contacto, gestión del PIN y atomicidad de la
escritura.

Una de las comprobaciones verifica que el bloqueo de escritura se adquiere antes
de la primera lectura. Sin esa garantía, dos envíos simultáneos leen el mismo
saldo, ambos superan la comprobación y una actualización sobrescribe a la otra:
en la reproducción del caso, dos envíos de 600 € sobre un saldo de 1.000 €
resultaban ambos aprobados.

\subsection{Limitaciones conocidas}

La evaluación permite delimitar con precisión lo que el sistema no resuelve:

\begin{itemize}
    \item \textbf{Resolución de referencias entre turnos.} Ante la secuencia
    ``¿cuánto he gastado en supermercado este año?'', ``¿y en restaurantes?'',
    ``¿y el año pasado?'', el modelo resuelve la tercera pregunta contra la
    primera y responde sobre supermercado. Es un comportamiento sistemático, no
    ocasional.

    \item \textbf{Previsión de meses futuros.} La herramienta proyecta
    únicamente el mes en curso. Ante una pregunta sobre el mes siguiente, el
    modelo responde con la cifra del mes actual sin advertir la diferencia. Se
    probaron dos correcciones —una regla en el \textit{prompt} y un campo
    explícito en la respuesta de la herramienta— y ninguna resultó efectiva.

    \item \textbf{Preguntas sin periodo explícito y consultas sobre terceros.}
    Se documentan como limitaciones asumidas.
\end{itemize}

Al margen de lo que mide el banco de pruebas, el sistema tiene otros límites
que conviene dejar por escrito, porque marcan la distancia entre esta
demostración y un producto.

\subsubsection{Alcance funcional}

\begin{itemize}
    \item \textbf{Un único titular y una única cuenta.} La tabla
    \texttt{cliente} contiene una sola fila y las consultas de saldo se hacen
    contra \texttt{WHERE id = 1}. La tabla \texttt{movimientos} no tiene columna
    de cuenta, así que no hay forma de distinguir entre varias.

    \item \textbf{Catálogo de operaciones reducido.} Se pueden consultar el
    saldo y los movimientos y enviar un Bizum. No hay transferencias, tarjetas,
    domiciliaciones ni cancelación de una operación ya realizada.

    \item \textbf{Histórico relativo a la fecha de generación.} Los movimientos
    se crean a partir del día en que se ejecuta \texttt{seed.py}. Si cambia el
    mes y no se regenera la base de datos, las preguntas sobre ``este mes''
    devuelven resultados vacíos.

    \item \textbf{Solo español.} El \textit{system prompt}, el reconocimiento de
    voz (\texttt{es-ES}) y los formatos de cifra y fecha asumen ese idioma.
\end{itemize}

\subsubsection{Límites técnicos y de operación}

\begin{itemize}
    \item \textbf{Memoria de la conversación acotada.} La ventana es de 8192
    \textit{tokens} y se conservan 24 mensajes. Una conversación larga pierde el
    principio, y el historial vive en memoria del proceso: al cerrarse la
    conexión por inactividad, o al reiniciar el servidor, desaparece.

    \item \textbf{Atención en serie.} Hay un solo proceso y una sola instancia
    del modelo en la tarjeta gráfica, que no admite dos peticiones a la vez. Con
    varios usuarios simultáneos las peticiones se encolan y la latencia crece de
    forma proporcional.

    \item \textbf{SQLite admite un único escritor.} La transacción
    \texttt{BEGIN IMMEDIATE} garantiza la integridad de los envíos, pero lo hace
    serializándolos.

    \item \textbf{Las latencias dependen del equipo.} Las cifras de esta memoria
    se midieron sobre un portátil concreto y conectado a la corriente; en otro
    equipo, o con la batería limitando la tarjeta gráfica, cambian de forma
    apreciable.

    \item \textbf{La interfaz necesita conexión.} El modelo y los datos son
    locales, pero el navegador descarga de la red la biblioteca de gráficos
    (Vega-Lite), las tipografías y los iconos de los comercios. Sin internet el
    asistente responde y habla igual, pero los gráficos no llegan a dibujarse.
\end{itemize}

\subsubsection{Límites de la demostración en materia de seguridad}

\begin{itemize}
    \item \textbf{El acceso no es autenticación real.} La pantalla de inicio se
    valida en el navegador y el punto de entrada \texttt{/ws} no comprueba
    ninguna credencial: quien conozca la dirección puede abrir una conversación.

    \item \textbf{El PIN está fijado en el código} (\texttt{config.py}), sin
    cifrado ni almacenamiento seguro. Su papel aquí es demostrar que la clave
    nunca llega al modelo, no proteger un activo real.

    \item \textbf{No hay registro de auditoría.} Las operaciones modifican el
    saldo pero no dejan traza de quién las pidió ni desde dónde.
\end{itemize}

\subsubsection{Límites de la propia evaluación}

\begin{itemize}
    \item \textbf{El resultado varía entre repeticiones} en uno o dos casos
    aunque no cambie nada del sistema. Por eso las cifras se obtienen siempre
    sobre tres repeticiones, y una diferencia de un caso no se interpreta como
    mejora ni como regresión.

    \item \textbf{El banco tiene 44 preguntas.} Cubre los tipos de consulta del
    enunciado y los añadidos durante el desarrollo, pero no es exhaustivo: hay
    formas de preguntar que no están representadas.

    \item \textbf{No se prueba la interfaz de forma automática.} Las pruebas
    llegan hasta los eventos que emite el servidor; lo que ocurre después en el
    navegador se ha comprobado a mano.
\end{itemize}

Delimitar estas limitaciones con precisión resulta más útil que un resultado
global elevado, ya que indica en qué condiciones conviene confiar en el sistema
y en cuáles no.
